\begin{proposition}[Cover-compatible unit locality]\label{prop:coverlocality}
Let $\pi:\widetilde G\to G$ be a graph covering used in a specified carrier or address registration. Then each incident edge at a selected lift projects bijectively onto the incident edges at its image, and each base walk has a unique lift from that selected start, preserving its length. If the admitted event grammar respects this correspondence, retains the information required for recovery, and gives each primitive field handoff address distance at most one, the unit causal bound of Proposition~22.2 holds in each description. A value $\cth=1$ additionally requires a saturating admitted depth-one field handoff in each domain being compared.
\end{proposition}
\begin{proof}
The neighborhood bijection is the defining covering condition. Induction gives the unique lift and equal walk length. Once the separate admission and field-distance hypotheses hold, the triangle-inequality proof of Proposition~22.2 applies in either description; an admitted adjacent depth-one event attains ratio one. A graph edge alone proves neither admission nor recovery nor saturation.
\end{proof}
The global multiplicity of a two-cover is two, while its local edge multiplicity at a selected start is one. This separates covering multiplicity from elementary walk length. It does not identify an internal $G_{60}$ edge with a field-address acquisition or equate endpoint distances for arbitrary projected histories.
